\section{B. Other logarithmic transformations and finite covers}

\subsection{B.1. Comparison under logarithmic transformation}

\begin{claim}{B1.0}{Proved}{The marked geometric dictionary on the first page}
The lattice matrix, meridian orientations, line bundle and parameter
comparison stated on the first page identify the two marked smooth
families. The ambiguity in $u$ is exactly multiplication by the
nonzero constant specifying Engel's reference linearization.
\end{claim}
\begin{proof}
First quotient the period presentation by its last three columns,
using $z=\zeta_1$ and $v=e^{2\pi i\zeta_2}$. The resulting principal
$\C^*$-bundle over $\C/(\Z+\tau\Z)$ has factors
\[
 (z,v)\mapsto(z+1,v),\qquad
 (z,v)\mapsto(z+\tau,e^{2\pi i\mu}v).
\]
For an odd theta function $\vartheta$, the meromorphic function
$\vartheta(z-\mu)/\vartheta(z)$ has precisely these factors and
divisor $(\mu)-(0)$. Thus this is the complement of the zero section
of $L_\mu=\OO(P-O)$, with $P=\mu$ on the elliptic curve. The first
period column acts by
$(z,v)\mapsto(z+6\mu,e^{2\pi i\beta}v)$; its inverse covers $-6P$,
as does Engel's contracting generator. Replacing $\beta^{\rm part}$
by $\beta^{\rm part}+c$ multiplies this inverse lift by
$e^{-2\pi i c}$. A change of reference lift multiplies its scale by
a single constant $\lambda_0\in\C^*$, proving the asserted parameter
comparison and naming its entire ambiguity.

For the integral markings we apply the displayed matrix $Q$ to
column lattice vectors. Exact integer multiplication gives
$\det Q=-1$, $T_0Q=QA_1^{-1}$ and $T_1Q=QA_2^{-1}$, as recorded
with the matrices in \texttt{b\_dictionary.py}. These inverses are
required because the source meridians have opposite orientations.
The product of the meridians gives
$QM_0Q^{-1}=T_0^{-1}T_\infty^{-1}T_0$ with the stated cusp path.
Finally $A_jv_j=v_j$, so inversion of an affine deck generator with
translation $v_j/m_j$ changes that translation to $-v_j/m_j$.
Applying $Q$ gives $a_0=-Qv_1/3$ and $a_1=-Qv_2/4$.
This calculation identifies marked lattices and affine deck actions;
it does not discard integral collar translations as if they were
equal to zero.
\end{proof}

\begin{claim}{B1.1}{Proved}{The relative cohomology comparison}
Fix the period family and cusp filling. Let $X'$ be obtained by any
admissible translation twists and multiplicities $m_1=3k$, $m_2=4l$,
where $k,l$ are positive integers. Write $D'=W+S'_1+S'_2$ and
\[
E'=\Omega^1_{X'}(\log D')/f'^*\Omega^1_B(\log\Sigma).
\]
The identity on the common smooth family induces isomorphisms on $B$
\[
R^qf_*\OO_X\simeq R^qf'_*\OO_{X'},\qquad
R^qf_*E\simeq R^qf'_*E'\quad(0\le q\le2).
\]
These isomorphisms identify the logarithmic connecting maps $\kappa_q$.
They also hold when the cusp collar is changed by an admissible translation.
\end{claim}
\begin{proof}
We construct the comparison at every point. Over $B\setminus\Sigma$
the two families are identified, and a translation acts trivially on
$H^q(F,\OO_F)$ and $H^q(F,\Omega^1_F)$ for a complex torus $F$.
For the second assertion, translation fixes the invariant holomorphic
one-forms, and the first follows either from invariant Dolbeault forms
or the exterior algebra $H^\bullet(F,\OO_F)=\bigwedge^\bullet H^1(F,\OO_F)$.
These are statements about the induced maps on cohomology, including
their holomorphic variation.

At a point with original order $d=3$ or $4$, let
$\widetilde f:T\to\Delta_z$ be the smooth proper torus model with
coarse coordinate $q=z^d$. The new model is its pullback by $z=s^e$,
where $m=de$, followed by the free cyclic quotient of order $m$.
Smooth proper base change identifies the higher direct images upstairs
with the pullbacks of those for $T$. The relative logarithmic cotangent
bundle upstairs is precisely $\Omega^1_{T/\Delta_z}$ pulled back:
the base logarithmic direction is killed in the quotient defining $E'$.
The derivative of a translation contributes only a multiple of $ds$;
it consequently acts on this quotient as ordinary translation on relative forms.

Average over the cyclic action and choose holomorphic eigenframes.
If a frame for $H^q(\OO)$ or $H^q(\Omega^1)$ has character
$\zeta_d^{-a}$, $0\le a<d$, its original invariant extension is
$z^a\OO_{\Delta_q}$. On the order-$m$ cover its character is
$\zeta_m^{-ea}$, and its invariant extension is
$s^{ea}\OO_{\Delta_q}=z^a\OO_{\Delta_q}$.
Taking invariants is exact in characteristic zero; hence this equality
is an equality of the extensions inside the common punctured-disc
bundle, not only an equality of their ranks or degrees. A translation
twist induces the identity on the cohomology just used and changes none
of these lattices. This proves the comparison at both finite points,
including the full nonreduced fibres.

At the cusp the relative logarithmic forms have the global toric frame
$d\log x_1,d\log x_2$, modulo $d\log q$.
A toric gluing has the form
$x_i\mapsto c_i(q)q^{b_i}x_i$ and therefore changes $d\log x_i$ by
$(q c_i'/c_i+b_i)d\log q$. It induces the identity on the relative
logarithmic quotient. The structure-sheaf comparison is also the identity
on the normalization cochain complex. This verifies the cusp extension
and the assertion about translated collars.

The logarithmic cotangent exact sequences are isomorphic on the smooth
base, so their connecting maps agree there under these comparisons.
Their source and target sheaves are locally free by the local calculations
in Part A and the eigenframe calculation above. A holomorphic map between
vector bundles on a reduced curve which is zero on a dense open set is
zero everywhere. Thus the connecting maps agree globally.
\end{proof}

\begin{claim}{B1.2}{Proved}{All admissible twists have the same Hodge numbers}
For every family in B1.1 whose cusp is the one-component hexagonal filling,
the direct images of the absolute cotangent sheaves are
\[
\begin{array}{c|ccc}
 &R^0f'_*&R^1f'_*&R^2f'_*\\\hline
\Omega^1_{X'}&\OO(-2)&\OO\oplus\OO(-1)^2&\OO\oplus\OO(-1)\\
\Omega^2_{X'}&\OO(-2)\oplus\OO(-1)&\OO(-2)\oplus\OO(-1)^2&\OO
\end{array}
\]
and $R^qf'_*\OO_{X'}=(\OO,\OO\oplus\OO(-1),\OO(-1))$.
In particular every admissible $(3,4)$ member with
$|12l_0-4l_1-3l_2|=1$ has the Hodge diamond of Part A.
\end{claim}
\begin{proof}
B1.1 identifies the logarithmic direct images and their connecting maps.
In Part A the cokernel of $\kappa_0$ is
$\C_{p_1}\oplus\C_{p_2}$, and $\kappa_1$ is surjective with kernel
$\OO\oplus\OO(-1)^2$. These statements therefore hold for $X'$.
The residue sheaf of its normalized divisor has higher direct images
\[
R^0f'_*\OO_{\widetilde D'}=\C_{p_0}\oplus\C_{p_1}\oplus\C_{p_2},
\quad R^1f'_*\OO_{\widetilde D'}=\C_{p_1}\oplus\C_{p_2},
\quad R^2=0.
\]
Indeed the finite reductions are bielliptic, and the cusp normalization
is rational. The base logarithmic residue is evaluation with weights
$(1,m_1,m_2)$; its kernel is $\OO(1)(-\Sigma)=\OO(-2)$.
On $H^1$ the residue of the invariant class $\gamma^{0,1}$ is
$m_j\gamma^{0,1}|_{S'_j}\ne0$. The other eigenline has zero residue.
The local rank-one image of $\kappa_0$ is its kernel, as follows either
from the residue diagram of Part A or from the eigenframe calculation
in B1.1. Thus the two finite skyscrapers cancel exactly. This proves
the first row, with no surviving torsion.

The proper flat map is Gorenstein of relative dimension two. Relative
duality for $F=\Omega^1_{X'}$, using
$F^*\otimes\omega_{X'/B}=\Omega^2_{X'}\otimes f'^*\OO(2)$,
gives the second row because all sheaves in the first row are locally
free. Coherent Leray on $\Pone$ gives every Hodge number. In particular
$h^{p,0}=0$ for $p>0$ follows from these direct images and duality,
so no separate invariance assumption about absolute forms is needed.
\end{proof}

\subsection{B.2. Higher multiplicities and the anticanonical ring}

\begin{claim}{B2.1}{Proved}{Canonical bundles for all coprime multiplicities}
Let $X=X_{3k,4l}(u)$ with admissible translation data as in
\citet[Theorem 1.1 and Proposition 2.3]{ViaclovskyLog2026},
$\gcd(3k,4l)=1$, and any nonzero $u$ in its construction domain.
Then
\[
K_X=f^*\OO_B(-1)\otimes
\OO_X\bigl((k-1)S_1+(3l-1)S_2\bigr).
\]
Moreover $(h^{0,q})_{q=0}^3=(1,1,0,0)$, $h^{p,0}=0$ for $p>0$,
and $(h^{1,1},h^{1,2})=(2,1)$ for every such pair.
Thus these Hodge numbers do not distinguish the multiplicity families.
\end{claim}
\begin{proof}
The cited proposition supplies exactly the free cyclic quotients used
in B1.1; its Theorem 1.1 supplies the stated sphere topology for its
admissible choices. B1.2 computes the cohomology, with no use of that
topological conclusion.
For the canonical bundle we check the natural evaluation
$f^*f_*\omega_{X/B}\to\omega_{X/B}$.
The invariant relative determinant has weights $a_1=2k$ and $a_2=l$
under the clockwise generators of orders $3k$ and $4l$. The pullback
of a canonical form is $F(s)\,ds\wedge e$, where the character of
$ds\wedge e$ is $\zeta_m^{1+a}$. Therefore its invariant coefficients
are $s^{m-1-a}$ times an invariant function and a nowhere vanishing
unit. Since the quotient of total spaces is etale, the evaluation
vanishes to order $m-1-a$ on the reduced fibre downstairs. These orders
are $k-1$ and $3l-1$. At the cusp the form
$d\log x_1\wedge d\log x_2\wedge dq$ is a nowhere vanishing canonical
frame in every smooth toric chart, so no cusp component occurs.

By B1.1 and relative duality,
$f_*\omega_{X/B}=(R^2f_*\OO_X)^*=\OO(1)$.
The evaluation divisor just computed therefore gives
$\omega_{X/B}=f^*\OO(1)\otimes
\OO((k-1)S_1+(3l-1)S_2)$. Tensoring by $f^*K_B=f^*\OO(-2)$
proves the formula.
\end{proof}

\begin{claim}{B2.2}{Proved}{The complete graded anticanonical ring}
Put $m_1=3k$, $m_2=4l$, and $d=4l+3k-kl$.
Choose integers $r,s$ with $m_2r+m_1s=1$ and set
$\mathcal L=\OO_X(rS_1+sS_2)$.
There is a graded algebra isomorphism
\[
\bigoplus_{n\ge0}H^0(X,-nK_X)
\simeq\bigoplus_{n\ge0}\C[x,y]_{dn},
\qquad\deg x=m_2,\quad\deg y=m_1,
\]
where the right-hand summand is assigned degree $n$.
For negative weights the polynomial space is zero. Explicitly
\[
h^0(X,-nK_X)=
\max\left\{1+n-\left\lceil\frac{n(k-1)}{3k}\right\rceil
-\left\lceil\frac{n(3l-1)}{4l}\right\rceil,\,0\right\}.
\]
The integer $d$ is nonzero for admissible pairs. If $d<0$ the
anticanonical ring is $\C$; if $d>0$ it is the indicated weighted
Veronese ring. For $(k,l)=(1,1)$ it is
$\C[U,V]$, $\deg U=1$, $\deg V=2$.
\end{claim}
\begin{proof}
The divisor of $f^*(t/(t-1))$ is $m_1S_1-m_2S_2$.
It follows, retaining this rational function as the comparison, that
$\mathcal L^{m_2}\simeq\OO(S_1)$,
$\mathcal L^{m_1}\simeq\OO(S_2)$ and
$\mathcal L^{m_1m_2}\simeq f^*\OO(1)$.
B2.1 then gives $-K_X=\mathcal L^d$.

We first establish the required section algebra for every integer
power of $\mathcal L$. If a meromorphic function represents a section
of $\OO(aS_1+bS_2)$, it is holomorphic on every regular compact fibre
and is constant there. It descends to the base on the smooth locus.
On a cyclic cover $q=s^m$, the finite pole bound gives a Laurent
expansion bounded below, and invariance leaves only powers divisible
by $m$. Thus descent extends meromorphically across the finite point,
with the exact valuation inequality $m\operatorname{ord}_q(h)+a\ge0$.
The identical argument at the reduced cusp has $m=1$.
Conversely each base function satisfying these bounds pulls back to
an allowed section. Consequently, as subsheaves of meromorphic functions,
\[
f_*\OO(aS_1+bS_2)=
\OO_B\left(\left\lfloor a/m_1\right\rfloor p_1+
\left\lfloor b/m_2\right\rfloor p_2\right).
\]
This proves the displayed dimension formula using B2.1 and
$h^0(\Pone,\OO(e))=\max\{e+1,0\}$.

For the algebra, let $x,y$ be the canonical sections of
$\OO(S_1),\OO(S_2)$ under the chosen identifications with powers of
$\mathcal L$. Their powers $x^{m_1},y^{m_2}$ pull back two independent
sections of $\OO_B(1)$, with zeros at the two distinct points.
For any integer $j$, choose the unique residues
$0\le a_0<m_1$, $0\le b_0<m_2$ for which
$j-m_2a_0-m_1b_0=e m_1m_2$ is integral.
The valuation formula identifies $H^0(\mathcal L^j)$ with
$x^{a_0}y^{b_0}H^0(\OO_B(e))$. When $e\ge0$ a basis is
\[
x^{a_0+m_1i}y^{b_0+m_2(e-i)},\qquad0\le i\le e;
\]
when $e<0$ there are no sections. These are exactly the monomials of
weighted degree $j$. Their multiplication is ordinary polynomial
multiplication, proving the algebra assertion in every degree.
When $k=l=1$, the weight-$6n$ monomials are generated by $y^2$ and
$x^3$, of regraded degrees $1,2$.
Finally, $d=0$ would imply $k\mid4l$ and $l\mid3k$; coprimality
forces $k=l=1$, which instead gives $d=6$.
\end{proof}

For completeness the requested small pairs give the following exact data.
Here $r_1,r_2$ are the coefficients in B2.1 and $P_n^-=h^0(-nK_X)$.
\[
\begin{array}{c|c|c|c|c|ccc}
(k,l)&(m_1,m_2)&(r_1,r_2)&(\deg x,\deg y)&d&P^-_1&P^-_2&P^-_3\\\hline
(1,1)&(3,4)&(0,2)&(4,3)&6&1&2&2\\
(1,2)&(3,8)&(0,5)&(8,3)&9&1&1&2\\
(3,1)&(9,4)&(2,2)&(4,9)&10&0&1&1\\
(1,4)&(3,16)&(0,11)&(16,3)&15&1&1&1\\
(5,1)&(15,4)&(4,2)&(4,15)&14&0&1&1\\
(1,5)&(3,20)&(0,14)&(20,3)&18&1&1&1
\end{array}
\]
The program \texttt{b\_families.py} enumerates all coprime pairs with
$m_1m_2\le60$ and checks the floor formula against exact weighted
monomial counting for $0\le n\le12$. These checks accompany the
all-degree proof above.

\subsection{B.3. Finite fundamental groups and universal covers}

\begin{claim}{B3.1}{Proved}{Connected cyclic covers and the cusp comparison}
For an admissible $(3,4)$ member with
$n=|12l_0-4l_1-3l_2|>1$, let $\pi:Y\to X$ be the universal
cover and $g=f\pi$. It is a finite etale holomorphic map of degree $n$.
Every fibre of $g$ is connected; its finite reductions are bielliptic,
its multiplicities are $3,4$, and its cusp $W_n$ normalizes to a
disjoint union of $n$ copies of $\mathrm{dP}_6$.
The pullback maps induce isomorphisms
\[
R^qf_*\OO_X\simeq R^qg_*\OO_Y,\qquad
R^qf_*E_X\simeq R^qg_*E_Y.
\]
\end{claim}
\begin{proof}
We first derive the fundamental-group presentation from the fillings.
The smooth family over the punctured base has zero section, so its
fundamental group is $\Lambda\rtimes\langle x,y\rangle$, where
$x,y$ are the clockwise meridians and act by $A_1,A_2$.
On a finite filling the lifted affine generator has $m_j$th power
translation by $v_j$: its $m_j$ translations add to $v_j$ because
$A_jv_j=v_j$. Thus filling the normal disc imposes
$x^3=t_{v_1}$ and $y^4=t_{v_2}$, with the signs fixed by the
clockwise convention on the first page.

The toric covering space of the cusp tube is simply connected.
Indeed its dense torus has cocharacter lattice $\Z^3$; a loop may
be perturbed into that torus, and the normal discs of the three rays
$(0,1),(e_1,1),(e_2,1)$ kill a basis of this lattice.
This argument applies also to the infinite fan, since a compact loop
and a homotopy meet only finitely many of its locally finite charts.
The same conclusion holds for the small tube: if a loop in
$|t_c|<\epsilon$ bounds a disc in the toric space, the torus action
$(x_1,x_2,t_c)\mapsto(x_1,x_2,a t_c)$ for a sufficiently small
positive $a$ moves that compact disc into the tube, and the scaling
homotopy of its boundary stays in the tube.
The cusp quotient therefore has fundamental group
$\Lambda/\Lambda_{\rm tor}$, and its toric meridian bounds the
normal disc at fixed $x_1,x_2$. A translated collar makes this
meridian $t_\mu y^{-1}x^{-1}$, where $\gamma(\mu)=l_0$.
Van Kampen consequently kills $\Lambda_{\rm tor}$ and imposes
$xy=t_\mu$.

The normal closure of $\Lambda_{\rm tor}$ in the fibre lattice is
$\ker\gamma$: it contains $\hat w,\hat\delta$ and
$A_1\hat w+\hat w=\hat u$, and all these relations lie in
$\ker\gamma$. The remaining generator $c$, the image of
$\hat\gamma$, is central because $\gamma A_j=\gamma$.
The presentation is thus
$\langle c,x,y\mid c\text{ central},\ xy=c^{l_0},\
x^3=c^{l_1},\ y^4=c^{l_2}\rangle$.
Eliminating $y$ makes it abelian, with relation matrix
\[
 \begin{pmatrix}3&-l_1\\4&l_2-4l_0\end{pmatrix},
 \qquad\det=-p.
\]
Its first Smith divisor is one because its entries include three
and four; hence the group is $\Z/|p|$.
After imposing $c=1$ the quotient has $x^3=x^4=1$, so it is trivial.
Thus the image of the regular-fibre lattice, containing $c$, generates
the entire group $\Z/n$. The same is true of each reduced special
fibre. At a finite point the cyclic model is a smooth torus family
over a disc, with a free affine finite-group action. An equivariant
smooth trivialization, obtained by averaging a horizontal connection,
contracts the disc to its centre and descends to a deformation
retraction of the local total space onto $S_j$. Thus the image of
the regular-fibre lattice in the local fundamental group is contained
in the image of $\pi_1(S_j)$; the same holds after mapping to
$\pi_1(X)$. At the cusp $\hat\gamma$ survives in
$\Lambda/\Lambda_{\rm tor}$. The explicit normal-crossing contraction
of the tube onto $W$ given in B3.3 identifies their fundamental groups,
so the same generator is in the image of $\pi_1(W)$.
The covering-space criterion therefore makes the inverse image of
every fibre connected. Admissibility gives $3\nmid l_1$ and odd $l_2$,
so $\gcd(n,12)=1$. The etale cover preserves fibre multiplicities.
On a finite reduction it cannot kill the nontrivial order-$3$ or
order-$4$ linear holonomy, since the covering degree is prime to it;
the resulting finite torus quotient is again bielliptic.

Normalization commutes with the finite etale base change $Y\to X$.
The pullback of the simply connected normalization $\mathrm{dP}_6$
is its trivial degree-$n$ cover. The double curves and triple points
also pull back etale, and their incidence maps are the lifts of the
original incidence maps. Consequently the dual complex is an actual
degree-$n$ covering of the original real two-torus, rather than merely
a cell complex with the same cell counts. It is connected because
$W_n$ is connected. The normalization resolution for $\OO_{W_n}$
has acyclic components, double curves and point strata. Its global
cochain complex is consequently the cellular cochain complex of this
dual torus, with $n,3n,2n$ cells. Thus its coherent cohomology is
$(1,2,1)$, and pullback from $W$ is an isomorphism over $\C$:
on a finite cover of real tori it is induced by a full-rank map of
their integral first cohomology lattices.

Modulo $d\log q$, the two toric forms $d\log x_i$ are invariant
under all cusp gluings as computed in B1.1. Hence
$E_X|_W=\OO_W^2$ and $E_Y|_{W_n}=\OO_{W_n}^2$,
and the same pullback argument applies to their cohomology.
Over a regular fibre, an isogeny of complex tori induces isomorphisms
on $H^q(\OO)$ and $H^q(\Omega^1)$ over $\C$.
At a finite fibre, make the cyclic base change. The map is an isogeny
of the smooth torus families and induces isomorphisms of their
cohomology bundles. Taking the same cyclic invariant eigenspaces
descends these isomorphisms, including all nilpotent layers of the
scheme fibre. Grauert base change at the cusp, and this explicit
calculation at the finite points, prove the asserted bundle maps are
isomorphisms on all of $B$.
\end{proof}

\begin{claim}{B3.2}{Proved}{Cotangent direct images and Hodge numbers of the covers}
Write $T_n=\C_{p_0}^{\,n-1}$. Then
\[
\begin{array}{c|ccc}
 &R^0g_*&R^1g_*&R^2g_*\\\hline
\Omega^1_Y&\OO(-2)&\OO\oplus\OO(-1)^2\oplus T_n&\OO\oplus\OO(-1)\\
\Omega^2_Y&\OO(-2)\oplus\OO(-1)&\OO(-2)\oplus\OO(-1)^2&\OO\oplus T_n
\end{array}
\]
The splittings involving $T_n$ are not canonical.
The only nonzero Hodge numbers of $Y$ are
\[
h^{0,0}=h^{0,1}=h^{3,2}=h^{3,3}=1,\quad
h^{1,2}=h^{2,1}=1,\quad h^{1,1}=h^{2,2}=n+1.
\]
The original finite-fundamental-group threefold $X$ itself has the
same Hodge diamond as the simply connected $(3,4)$ member.
\end{claim}
\begin{proof}
B3.1 identifies the logarithmic cotangent sequence on cohomology,
and therefore its connecting maps and direct images, with those of
the original family. The normalized residue sheaf differs only at
the cusp: its degree-zero cohomology there is now $\C^n$.
The residue of a base logarithmic form is the diagonal vector
$(1,\ldots,1)$, since each cusp component has multiplicity one.
Its cokernel is $T_n$. The degree-one residues at the finite fibres
cancel the same two skyscrapers as in B1.2. The long exact sequence
therefore gives
\[
0\longrightarrow T_n\longrightarrow R^1g_*\Omega^1_Y
\longrightarrow\OO\oplus\OO(-1)^2\longrightarrow0
\]
and the remaining entries of the first row. The extension splits
because $\operatorname{Ext}^1(V,T_n)=H^1(V^*\otimes T_n)=0$
for a vector bundle $V$ on $B$.

For clarity relative duality must now be used in its derived form.
For a sheaf $F$ locally free on the total space, it gives
\[
Rg_*(F^*\otimes\omega_{Y/B})
\simeq R\mathcal Hom_B(Rg_*F,\OO_B)[-2].
\]
For $F=\Omega^1_Y$, the torsion in $R^1g_*F$ contributes
$\mathcal Ext^1_B(T_n,\OO_B)$ in degree two of the right side.
This is a length-$(n-1)$ skyscraper at $p_0$. After tensoring by
$\OO(-2)$ it still has that length. The locally free parts dualize
in complementary degrees, giving exactly the second row; the
extension by the locally free quotient splits for the same reason
as above. Coherent Leray computes the displayed Hodge numbers.
For $X$ with one-component cusp, B1.2 already applies independently
of its fundamental group.
\end{proof}

\begin{claim}{B3.3}{Proved}{Canonical rings and rational topology of the covers}
For $X$ and $Y$ as above,
\[
K_Y=\pi^*K_X=\OO_Y(-2S'_2),\quad
\bigoplus_{m\ge0}H^0(Y,-mK_Y)=\C[U,V],\quad
\deg(U,V)=(1,2).
\]
The same ring holds for $X$.
Over $\C$, the Betti numbers of $Y$ are
$(1,0,n-1,0,n-1,0,1)$; its Fr\"olicher sequence degenerates at $E_2$.
For every nontrivial character of the deck group, the character
part of Dolbeault cohomology consists of one class in $(1,1)$ and
one in $(2,2)$ and no others.
\end{claim}
\begin{proof}
The canonical formula pulls back through an etale map, and
$4S'_2=g^*p_2$. B3.1 gives $g_*\OO_Y=\OO_B$.
The valuation argument in B2.2 therefore applies unchanged and proves
the ring statement with all multiplication maps.

The comparisons of B3.1 are pullback maps and identify the logarithmic
cohomology bundles with the invariant part; the deck group acts
trivially on those bundles. At the cusp its action on the $n$
normalization components is the regular permutation representation.
The quotient by the diagonal has each nontrivial character once.
Thus the first-row additional skyscraper in B3.2 has precisely these
characters, and duality gives the conjugate characters in its
second-row additional skyscraper. All other cohomology is invariant
and equal to that of $X$. Coherent Leray now proves the assertion
about the Dolbeault character spaces.

The de Rham complex and its Fr\"olicher sequence split under the
idempotents of the finite deck group. A nontrivial character has just
the two stated Dolbeault entries. No differential can connect them
or connect either to another nonzero entry, so both survive.
The invariant summand is the de Rham complex of $X$ by finite
etale descent and averaging. We compute its rational topology here
from the actual attachment maps; the integral homology assertion
of AC26 is not used as an input. In the remainder of this proof
all homology and cohomology coefficients are $\Q$.

Let $N$ be a small closed cusp tube, $M=\partial N$, and
$U=X\setminus\operatorname{int}N$. The space $U$ is covered by
two pieces retracting onto $S_1,S_2$, whose intersection retracts
onto a regular fibre $F$. The retractions are obtained from the
equivariant smooth torus trivializations used in B3.1. The maps
$F\to S_j$ are the central finite covering maps; translation of
their origins induces the identity on their rational cohomology.
Finite-cover transfer identifies their pullbacks with the invariant
subspaces for $A_j$. Put
\[
 \psi_1=2u+w+3\delta,\qquad \psi_2=u+w+2\delta,
 \qquad \eta=u\wedge w+6\gamma\wedge\delta.
\]
Direct substitution in $A_1,A_2$ shows that the invariant
degree-one spaces are $\langle\gamma,\psi_j\rangle$ and the
invariant degree-two spaces are
$\langle\gamma\wedge\psi_j,\eta\rangle$.
Their dimensions are two: the nontrivial eigenvalues of $A_j$ on
the complementary plane are the inverse primitive third or fourth
roots, so the exterior-square fixed space has exactly the two
displayed generators.
In the transfer-normalized dual homology bases, the two
Mayer--Vietoris maps are therefore the following actual matrices:
\[
 \alpha_1=\begin{pmatrix}
 1&0&0&0\\0&2&1&3\\-1&0&0&0\\0&-1&-1&-2
 \end{pmatrix},\qquad
 \alpha_2=\begin{pmatrix}
 2&1&3&0&0&0\\0&0&6&1&0&0\\
 -1&-1&-2&0&0&0\\0&0&-6&-1&0&0
 \end{pmatrix}.
\]
The second matrix uses the ordered wedge basis
$(\hat\gamma\hat u,\hat\gamma\hat w,\hat\gamma\hat\delta,
\hat u\hat w,\hat u\hat\delta,\hat w\hat\delta)$.
Both matrices have rank three, and
$\ker\alpha_1=\Q\nu_0$, with
$\nu_0=-\hat u-\hat w+\hat\delta$.
Consequently
\[
 H_1(U)=\Q,\qquad
 0\longrightarrow\operatorname{coker}\alpha_2=\Q
 \longrightarrow H_2(U)
 \xrightarrow{\partial}\Q\nu_0\longrightarrow0.       \tag{B.1}
\]
The generator $\Theta\in H^1(U)$ is normalized by
$\Theta|_F=\gamma$. The finite affine generators have
$3\bar g_1=l_1c$, $4\bar g_2=l_2c$ in rational homology, where
$c$ is the image of $\hat\gamma$ and $\Theta(c)=1$.
For the toric meridian $a_0=t_\mu y^{-1}x^{-1}$ this gives the
explicit collar value
\[
 \Theta(j_*a_0)=l_0-l_1/3-l_2/4=p/12.                 \tag{B.2}
\]

We justify the comparison with $N$ without assuming a collapse
theorem from AC26. In each normal-crossing chart
$t_c=z_1\cdots z_r$, $1\le r\le3$, choose the real vector field
whose $z_i$ component is
$-z_i|z_i|^{-2}/\sum_j|z_j|^{-2}$ off the divisor.
It satisfies $\dot t_c=-t_c$ and has norm at most
$C|t_c|^{1/r}$. A finite partition of unity preserves the first
identity and gives a global field with norm at most
$C'|t_c|^{1/3}$ after shrinking the tube. Extend it by zero on $W$.
Its flow exists for all positive time, and its tail length is
bounded by $3C'\epsilon^{1/3}e^{-T/3}$ after time $T$.
Thus the flow converges uniformly to $W$ and gives a continuous
strong deformation retraction fixing $W$. The same estimate proves
continuity at the divisor and at the endpoint of the retraction.

The constant-sheaf normalization resolution of $W$ has strata
$\mathrm{dP}_6$, three copies of $\Pone$, and two points.
Its row of degree-zero cohomology is
\[
 \Q\xrightarrow{0}\Q^3
 \xrightarrow{(a,b,c)\mapsto(a-b+c,a-b+c)}\Q^2.
\]
The strata have no degree-one cohomology. This resolution therefore
gives $H_1(N)=H_1(W)=\Q^2$, $H_4(N)=\Q[W]$ and $H_5(N)=0$;
it also gives $e(W)=6-3\cdot2+2=2$.
Poincare--Lefschetz duality for the compact smooth manifold with
boundary $(N,M)$ now shows that $H_1(M)\to H_1(N)$ and
$H_2(M)\to H_2(N)$ are surjective. For degree one the next
relative group is $H_1(N,M)=H^5(N)=0$.
For degree two the map to $H_2(N,M)=H^4(N)$ is the intersection
pairing with $[W]$, which is zero because $W$ is the divisor of
the global function $t_c$ on $N$. The Wang sequence for the
boundary mapping torus, with $\rank(M_0-I)=2$, gives
$\dim H_1(M)=3$. Its kernel to $H_1(N)$ is exactly
$\Q a_0$: this meridian is nonzero under projection to the base
circle and bounds its normal disc in $N$.

We exhibit two classes in
$K_2=\ker(H_2(M)\to H_2(N))$.
For $\nu=\hat w,\hat\delta$, let $\xi_\nu$ be the two-torus
obtained by translating the toric meridian around the angular
circle $S^1_\nu$. In the chart $t_c=z_0z_1z_2$ corresponding to
the rays $(0,1),(e_1,1),(e_2,1)$, these are the boundaries of
$D^2_{z_0}\times S^1_{z_i}$ with the other coordinate fixed
and the radii chosen so that $|t_c|=\epsilon$ on the boundary.
They therefore bound in $N$ with the literal toric meridian;
no choice of an unrecorded collar lift is being made.
Cutting the meridian into the two arcs of the cover defining
(B.1), the two cylinders have end circles differing by
$(I-A_1^{-1})\nu$. With this orientation of the Mayer--Vietoris
boundary, matrix multiplication gives
\[
 \partial j_*\xi_{\hat w}=-\nu_0,\qquad
 \partial j_*\xi_{\hat\delta}=0.                       \tag{B.3}
\]
The affine translation of either endpoint has no effect on its
homology class, which explains why this boundary formula is the
same for every translation twist.

Choose $\mathcal A\in H^2(U)$ restricting to $\eta$ on $F$;
it exists by Mayer--Vietoris because the two transfer-normalized
classes $\eta$ have equal restrictions. Fibrewise translation by
$S^1_{\hat\delta}$ defines a global smooth action on $U$:
$A_j\hat\delta=\hat\delta$, and all affine translations commute
with this circle. Slant product of its action map with the circle
defines $\iota:H^2(U)\to H^1(U)$. Naturality on $F$ gives
$\iota\mathcal A|_F=\iota_{\hat\delta}\eta=-6\gamma$.
The restriction $H^1(U)\to H^1(F)$ is injective by (B.1), so
$\iota\mathcal A=-6\Theta$. By the definition of slant product,
\[
 \langle\mathcal A,j_*\xi_{\hat\delta}\rangle
 =\langle\iota\mathcal A,j_*a_0\rangle=-p/2.           \tag{B.4}
\]
We may add the class pulled back from $(\Q\nu_0)^*$ under
the boundary map in (B.1) so that
$\langle\mathcal A,j_*\xi_{\hat w}\rangle=0$;
this does not change its restriction to $F$ or (B.4).
In the coordinates on $H_2(U)$ given by $\partial/\nu_0$ and
this normalized pairing, the two actual attachment cycles have
matrix
\[
 \bigl[j_*\xi_{\hat w}\ \ j_*\xi_{\hat\delta}\bigr]
 =\begin{pmatrix}-1&0\\0&-p/2\end{pmatrix}.             \tag{B.5}
\]
It is invertible over $\Q$ because admissibility implies
$p\ne0$. Thus the image of $K_2$ already fills $H_2(U)$.
Likewise (B.2) says the image of $\ker(H_1(M)\to H_1(N))$
is nonzero in $H_1(U)$. Using the surjectivity to $N$ just proved,
Mayer--Vietoris for $X=U\cup N$ now gives $H_1(X)=H_2(X)=0$.
Every other fibre has Euler number zero, so additivity over the
three special fibres and the smooth torus bundle gives $e(X)=2$.
Poincare duality then gives $b_4=b_5=0$ and $b_3=0$.
This proves the asserted rational homology of $X$ for all admissible
twists; the only changing determinant in (B.2)--(B.5) is a nonzero
constant multiple of $p$. The integer matrices and all exterior-power
identities were checked exactly in \texttt{code/b\_topology.py},
with output \texttt{outputs/b\_topology.txt}.

Part C's argument using this rational topology and the same Hodge
table now gives degeneration of the invariant summand at $E_2$.
Summing it with the nontrivial character spaces proves the Betti
numbers and the degeneration statement for $Y$. In particular
$e(Y)=2n$, as also follows from the covering degree.
\end{proof}
